% Modernised reading — fonds Alexandre Grothendieck, shelfmark 155, pages 1-10.
%
% This is an INTERPRETATION, not a transcription. Everything the critical
% apparatus carried moves into footnotes. In any dispute of substance the
% transcription governs: whoever cites Grothendieck must cite batch-01.fr.tex.
%
% Pass: Opus 5 (claude-opus-5), 2026-09-12 — first pass, unchecked against the
% pages by a human. The folder was read whole; it holds one batch, and of its
% ten leaves the five odd-numbered ones are his and are transcribed. The five
% even ones are unrelated paper and carry nothing of his.
%
% What the folder is. Five leaves, his own pagination 1 to 5, one argument:
% how to read off, from the algebra that a tensor category and a fibre functor
% over a FINITE EXTENSION k'/k produce, the fact that every object of the
% category has a dual. Page 1 states the dictionary (category + fibre functor
% <--> a commutative k'-algebra U with augmentation, plus an associative
% multiplication on A = U-dual); page 3 runs it on the smallest example, where
% U = k' tensor_k k' and A = End_k(k'), and then poses the folder's question;
% pages 5 to 9 attempt the answer and get as far as the antipode and the
% evaluation, breaking off mid-line on the last leaf.
%
% Notation fixed for the whole document. k a field, k'/k a finite extension.
% C is the tensor category, F : C --> Mod_f(k') the fibre functor (he writes F
% on page 1 and F' from page 5 on; one functor, written F' throughout here,
% and the reading says so). U is the commutative k'-algebra with augmentation
% epsilon, A = U-check = Hom_{k'}(U, k') the dual pro-object. A_g and A_d are
% A as a left and as a right module over itself (his "gauche"/"droite"). D is
% the duality of C, M-check^{(k')} the k'-linear dual. J = Ker(k' tensor_k k'
% --> k'), script-J the ideal it generates in A tensor-hat_k A, J-tilde the
% idealiser of that ideal.
%
% Three homonymies of the pages, fixed here and stated on the spine:
% — A is U-dual on page 1 and in the computation of page 3; but in the question
%   at the foot of page 3 the letter A carries an index and a gloss "= A
%   tensor_k k'", and there names the base change, an algebra over k' tensor_k
%   k'. The reading writes A_{k'} for that one.
% — D is the duality of the category on pages 5 and 7, and D_{k'}(A) in the
%   margin of page 5 is the k'-linear dual of A.
% — the two J-glyphs of page 1, which the leaf separates only by a tilde.
%
% Substantive departures from the pages, each footnoted where it occurs:
% — page 1, article 2°), qualifies the structure carried by A with a word the
%   transcription flags as doubtful and reads "commutative et associative".
%   Commutativity is excluded by the folder's own example on page 3, where
%   A = End_k(k') is a matrix algebra; the reading says "associative, d'unite
%   epsilon" and footnotes the page.
% — page 1 writes the comultiplication as A --> A tensor-hat_{k'} A and then
%   says its values lie in J-tilde/script-J. The second statement corrects the
%   first: A tensor-hat_{k'} A is not a ring, and the reading says why.
% — page 5 strikes "bimodule" and asks for a module structure on U over
%   A tensor A; a bimodule is a module over A tensor A-opposite, which is what
%   the next line of the same page (End_A(U) = A-opposite) in fact produces.
% — page 3's condition that the base-change functor be faithful does not hold
%   for an arbitrary module over k' tensor_k k'; the reading states what is
%   true and says what the condition must therefore be asking.
%
% Modern names given here and footnoted as ours: catégorie tannakienne, non
% neutre; foncteur fibre; groupoïde affine agissant transitivement sur Spec k';
% algébroïde de Hopf; antipode; produit de Takeuchi; objet dualisant; descente
% galoisienne; équivalence de Morita; Skolem-Noether. The leaves carry none of
% these words.
%
% What the folder announces and never establishes: the intrinsic conditions on
% C it sets out to find; the answer to its own "? ? ?" (does rigidity plus
% k = End(1_C) suffice); the verification of the last compatibility, begun on
% page 9 and abandoned mid-line; and the converse direction — that an antipode
% gives back rigidity — which the pages never state.
\documentclass[11pt,a4paper]{article}
% Le préambule commun à toutes les transcriptions du fonds Grothendieck.
%
% Il tient en une idée : **l'incertitude se compose, elle ne se résout pas.**
% Une transcription qui lisse un mot douteux a détruit la seule chose pour
% laquelle on la faisait — un lecteur doit pouvoir savoir, à chaque endroit,
% ce qui était sur le feuillet et ce qui a été ajouté. Les six macros qui
% suivent sont donc l'appareil critique, et non de la décoration.
%
% Les mêmes commandes sont reconnues par `scripts/render.mjs`, qui produit la
% vue de lecture du site. Une macro ajoutée ici doit l'être aussi là-bas,
% faute de quoi le rendu échouera bruyamment — ce qui est voulu : un rendu
% silencieusement faux est pire qu'un rendu absent.

\ProvidesPackage{grothendieck}[2026/08/08 Transcriptions du fonds Grothendieck]

\usepackage{amsmath,amssymb,amsthm}
\usepackage{mathrsfs}
\usepackage{stmaryrd}
% Les diagrammes commutatifs sont partout dans ce fonds : c'est la langue
% même de la géométrie algébrique de Grothendieck, et une transcription qui
% les rendrait en prose ne transcrirait rien.
\usepackage{tikz-cd}
% Français par défaut — c'est la langue des feuillets — mais l'anglais est
% chargé aussi : la traduction et le résumé sont des documents anglais, et la
% typographie française (espaces avant les deux-points) y serait une faute.
% Ces documents basculent par \selectlanguage{english} en tête de corps.
\usepackage[english,french]{babel}
\usepackage{fontspec}
\usepackage{geometry}
\geometry{a4paper,margin=25mm,marginparwidth=28mm}
\usepackage{marginnote}
\usepackage{xcolor}
% ulem, not soul. `soul`'s \st and \ul are built on a character-by-character
% scan that XeLaTeX's font machinery breaks: the document fails with an
% undefined control sequence rather than degrading. ulem does the same two jobs
% and is Unicode-engine safe. `normalem` stops it hijacking \emph, which the
% transcriptions use for Grothendieck's own emphasis.
\usepackage[normalem]{ulem}
\usepackage{hyperref}
\hypersetup{colorlinks=true,linkcolor=black,urlcolor=black,pdfborder={0 0 0}}

% --- Métadonnées du lot -----------------------------------------------------
% Reprises telles quelles par le script de rendu, qui les affiche en tête de
% la vue de lecture. Elles fixent ce que le fichier prétend couvrir : sans
% elles, une transcription flotte sans point d'attache dans un fonds de seize
% mille feuillets.

\newcommand{\folder}[1]{\gdef\grothFolder{#1}}
\newcommand{\batch}[1]{\gdef\grothBatch{#1}}
\newcommand{\leaves}[2]{\gdef\grothFirst{#1}\gdef\grothLast{#2}}
\newcommand{\foldertitle}[1]{\gdef\grothFoldertitle{#1}}
\gdef\grothFolder{?}\gdef\grothBatch{?}\gdef\grothFirst{?}\gdef\grothLast{?}\gdef\grothFoldertitle{}

% --- Le résumé --------------------------------------------------------------
%
% Il ouvre la lecture modernisée, sous le titre « Résumé », et il est composé
% en petit. Ce n'est pas un ornement : c'est le seul passage du document qui
% ne soit pas des mathématiques, et un lecteur doit pouvoir le reconnaître
% avant de l'avoir lu — puis le sauter s'il connaît déjà le sujet, ou s'y
% arrêter s'il ne le connaît pas. Le filet qui le suit marque la reprise du
% corps.

\newenvironment{resume}
  {\small\setlength{\parskip}{0.45em}}
  {\par\medskip\hrule\medskip}

% --- Mots-clés ---------------------------------------------------------------
%
% En anglais, en clôture du résumé de la lecture modernisée : le vocabulaire
% moderne sous lequel chercher ce que les feuillets construisent. Le manifeste
% du site (`npm run manifest`) les extrait pour étiqueter la cote entière —
% cette ligne est la source unique des tags, il n'y a pas de fichier à côté.
\newcommand{\keywords}[1]{\par\smallskip\noindent
  {\footnotesize\textsc{Keywords} --- \textit{#1}}\par}

% --- L'appareil critique ----------------------------------------------------

% Le numéro de feuillet, en marge. C'est l'ancre qui permet de régler un
% désaccord sur une formule en regardant *un* feuillet plutôt qu'une chemise
% de six cents. Le site s'en sert aussi pour tourner les pages du fac-similé
% quand on fait défiler la transcription.
\newcommand{\leaf}[1]{%
  \marginnote{\footnotesize\color{gray!70}\textsf{#1}}%
  \label{leaf:#1}\ignorespaces}

% Illisible : on ne devine pas. Un mot inventé qui se lit comme les autres est
% le pire résultat possible.
\newcommand{\ill}{\textcolor{red!60!black}{[\,\ldots\,]}}

% Lecture douteuse : le mot est donné, et signalé comme incertain.
\newcommand{\uncertain}[1]{\uline{#1}}

% Ajout de l'éditeur — une lettre manquante, un mot sous-entendu.
\newcommand{\add}[1]{\textcolor{blue!50!black}{[#1]}}

% Biffé par Grothendieck lui-même. Ce qu'il a rayé fait partie du document :
% une rature dit souvent où la pensée a changé de direction.
\newcommand{\struck}[1]{\sout{#1}}

% Note du transcripteur, sur l'état matériel du feuillet.
\newcommand{\note}[1]{\par\smallskip{\small\color{gray!60!black}\textit{[#1]}}\par\smallskip}

% Note marginale de Grothendieck, distincte du corps du texte.
\newcommand{\marginal}[1]{\par\smallskip\noindent\hspace{1em}%
  {\small\color{gray!60!black}#1}\par\smallskip}

% --- Titre ------------------------------------------------------------------

\AtBeginDocument{%
  \begin{center}
    {\large\bfseries Fonds Alexandre Grothendieck --- cote n\textsuperscript{o}~\grothFolder}\\[2pt]
    {\itshape\grothFoldertitle}\\[4pt]
    {\small feuillets \grothFirst--\grothLast\ (lot \grothBatch)}\\[2pt]
    {\footnotesize\color{gray!60!black}Transcription établie d'après le fac-similé numérisé
     par l'Université de Montpellier.\\
     Les lectures incertaines sont soulignées, les illisibles notées [\,\ldots\,],
     les ajouts entre crochets.}
  \end{center}
  \vspace{1em}}

\endinput


\folder{155}
\pages{1}{10}
\dating{[à partir de 1984]}
\watermark{Édition de démonstration\\interprétation personnelle de l'œuvre}
\foldertitle{[Rigidité de catégories tensorielles ?] : notes manuscrites (s.d.) — lecture modernisée du dossier entier}

\begin{document}

\section*{Résumé}

\begin{resume}
D'un groupe on connaît d'ordinaire les éléments. On peut aussi le connaître
autrement : par la liste de ses représentations linéaires — les manières de le
faire agir sur un espace vectoriel — et par la loi qui dit comment deux
représentations s'en combinent en une troisième. Le fait remarquable, et
ancien, est que cette seconde description suffit : le groupe se reconstruit à
partir d'elle. C'est le théorème de Tannaka pour les groupes compacts, et c'est
l'idée que ces feuillets appliquent ailleurs.

Ailleurs, parce qu'une difficulté s'ajoute dès qu'on travaille sur un corps
$k$ qui n'est pas algébriquement clos. Il arrive que les objets d'une famille
n'aient pas de description linéaire visible sur $k$, et qu'ils en acquièrent
une dès qu'on s'autorise une extension finie $k'$ de $k$ — comme un fibré qui
n'est trivial qu'au-dessus d'un revêtement, et pas au-dessus de l'espace
lui-même. On peut alors « voir » les objets, mais seulement après ce
changement de base, et ce qu'on reconstruit au bout du compte n'est plus un
groupe : c'est un objet à deux pattes, qui relie deux points de vue sur le
même objet plutôt que d'agir en un seul — un \emph{groupoïde}.

Le dossier est le brouillon du dictionnaire correspondant. D'un côté une
catégorie $C$ d'objets abstraits, munie d'un produit tensoriel et d'un
foncteur $F$ qui à chaque objet associe un espace vectoriel sur $k'$ ; de
l'autre une algèbre $A$ dont les modules sont exactement les objets de $C$, et
une comultiplication sur $A$ qui code le produit tensoriel. La première page
écrit ce dictionnaire en deux articles et deux conditions. La troisième le fait
tourner sur le plus petit exemple possible, celui où $C$ est la catégorie des
espaces vectoriels sur $k$ eux-mêmes et où $F$ est l'extension des scalaires :
on trouve $A = \operatorname{End}_{k}(k')$, l'algèbre de toutes les
applications $k$-linéaires de $k'$ dans lui-même — l'objet même dont la théorie
de la descente a besoin.

La question que le dossier pose alors est celle de son titre. Dans un groupe,
tout élément a un inverse. Du côté des représentations, cet inverse se voit
autrement : il dit que toute représentation a une \emph{duale}, et que
l'accouplement entre une représentation et sa duale est lui-même compatible à
tout le reste. Une catégorie où chaque objet a un dual de cette sorte est dite
\emph{rigide}. La question est : la rigidité de $C$, qui est une propriété
interne, purement catégorique, correspond à quelle propriété de l'algèbre $A$ ?
La réponse que les feuillets dégagent est que la rigidité fournit sur $A$ une
opération qui renverse les produits — $ab$ devient $\check{b}\check{a}$ — et
c'est exactement la trace, dans le dictionnaire, du passage à l'inverse dans un
groupe. Le nom moderne de cette opération est \emph{antipode}.

Deux choses méritent d'être signalées au lecteur qui suivra le détail. La
première est un accident propre à l'extension : $A$ n'est pas commutative, et
le corps $k'$ n'est pas central dedans — il y agit de deux manières, par la
gauche et par la droite, et les deux ne coïncident pas. Il s'ensuit que l'objet
qu'on voudrait écrire $A \otimes_{k'} A$ n'est pas une algèbre du tout, et
qu'il faut lui substituer un sous-quotient. La première page fabrique ce
sous-quotient à la main, sans lui donner de nom ; c'est ce que la théorie
appellera plus tard le produit de Takeuchi.

La seconde est la manière dont le dossier travaille. Il ne démontre presque
rien. Il pose un dictionnaire, le vérifie sur un exemple, énonce ce qu'il
faudrait, écrit trois points d'interrogation là où il ne sait pas, se donne un
programme à l'impératif — « Exprimer la rigidité de $C$ ! » — et s'arrête au
milieu d'une ligne quand la vérification qui reste devient une affaire de
calcul. Ce sont cinq feuillets d'orientation, pas un texte ; leur intérêt est
de montrer sous quelle forme une théorie se présente à quelqu'un avant d'être
écrite.

\keywords{tannakian category, non-neutral tannakian category, fibre functor, affine groupoid scheme, Hopf algebroid, antipode, rigidity, dualizable object, Takeuchi product, idealizer, Galois descent, Morita equivalence, pseudo-compact algebra, coefficient algebra, separable extension}
\end{resume}

\pagerange{1}{9}
\section*{Le fil du dossier, et les conventions}

Cinq feuillets, paginés de sa main de 1 à 5, qui sont les pages 1, 3, 5, 7 et 9
de l'inventaire\footnote{Les cinq feuillets intercalaires sont du papier
étranger — un contrat d'édition, trois pièces administratives de l'USTL et de
Paul Valéry, une circulaire — et ne portent rien de sa main. Leur seul usage
ici est de dater : ils sont de 1984, ce qui est la datation retenue par les
archivistes.}. L'argument ne s'interrompt pas d'un feuillet à l'autre et se
lit en cinq stations :

\begin{itemize}
\item \textbf{page 1} : le dictionnaire. Ce que c'est que se donner une
catégorie tensorielle sur $k$ et un foncteur fibre à valeurs dans les
$k'$-espaces vectoriels de dimension finie, en termes d'une algèbre.
\item \textbf{page 3, première moitié} : l'exemple de la descente,
$C = \operatorname{Mod}_{f}(k)$ et $F = -\otimes_{k} k'$, où le dictionnaire
donne $U = k' \otimes_{k} k'$ et $A = \operatorname{End}_{k}(k')$.
\item \textbf{page 3, seconde moitié} : la question du dossier, en deux
morceaux — une condition technique d'exactitude et de fidélité, et la question
de fond : la rigidité suffit-elle ?
\item \textbf{page 5} : la rigidité mise sous forme représentable. Le foncteur
de dualité est un $\operatorname{Hom}_{A}(-, I)$, l'objet dualisant $I$ est
$U$ lui-même, et la structure supplémentaire que cela exige sur $U$ équivaut à
un anti-isomorphisme de $A$.
\item \textbf{pages 7 et 9} : le calcul du dual, $F'(D(M)) = \check{M}^{(k')}$,
puis la compatibilité qui reste — que l'accouplement d'évaluation soit un
morphisme de $C$ — réduite au cas $M = A$, et abandonnée en cours de ligne.
\end{itemize}

\subsection*{Notation}

$k$ est un corps, $k'/k$ une extension finie. $C$ est une $\otimes$-catégorie
$k$-linéaire, associative, commutative et unifère\footnote{Le manuscrit abrège
cette liste d'adjectifs en trois lettres liées, que la transcription lit
« ACU » et signale comme incertaine ; l'abréviation n'est développée nulle part
dans le dossier. Elle est en usage dans la littérature tannakienne de ces
années-là, et c'est ce qui décide la lecture.}, et $F' : C \to
\operatorname{Mod}_{f}(k')$ le foncteur fibre : $k$-linéaire, exact, fidèle,
compatible au produit tensoriel\footnote{La page 1 l'appelle $F$ et les pages 5
à 9 l'appellent $F'$ ; c'est le même foncteur, et il est noté $F'$ partout
ici.}.

$U$ est une $k'$-algèbre commutative munie d'une augmentation $\varepsilon : U
\to k'$, et
\[
A \;=\; \check{U} \;=\; \operatorname{Hom}_{k'}(U, k')
\]
son dual, pris comme pro-objet de $\operatorname{Mod}_{f}(k')$ : $U$ est la
réunion filtrante de ses sous-espaces de dimension finie, $A$ la limite
projective de leurs duaux, c'est-à-dire une algèbre linéairement topologisée.
$A_{g}$ et $A_{d}$ désignent $A$ vue comme module à gauche, resp. à droite, sur
elle-même — les indices sont de lui. $D$ est la dualité de $C$ et
$\check{M}^{(k')} = \operatorname{Hom}_{k'}(M, k')$ le dual $k'$-linéaire.
Enfin
\[
J \;=\; \operatorname{Ker}(k' \otimes_{k} k' \to k'),
\qquad
\mathcal{J} \;=\; J\cdot(A \mathbin{\hat{\otimes}}_{k} A),
\qquad
\widetilde{\mathcal{J}} \;=\; \{\lambda \mid \lambda\mathcal{J} \subset
\mathcal{J}\}.
\]

\subsection*{Trois homonymies}

Elles rendent les feuillets illisibles pris ligne à ligne, et elles sont
levées une fois pour toutes ici.

\emph{La lettre $A$.} Elle désigne $\check{U}$ page 1 et dans le calcul de la
page 3. Mais dans la question posée au bas de la page 3 elle porte un indice et
une glose interlinéaire « $= A \otimes_{k} k'$ » : là, elle nomme le changement
de base, qui est une algèbre sur $k' \otimes_{k} k'$ et non sur $k'$. Cette
lecture l'écrit $A_{k'}$.

\emph{La lettre $D$.} C'est la dualité de la catégorie pages 5 et 7 ; le
$D_{k'}(A)$ de la note marginale de la page 5 est le dual $k'$-linéaire de $A$,
donc $U$.

\emph{Les deux $J$.} Le feuillet ne les distingue que par un tilde. $J$ est un
idéal de $k' \otimes_{k} k'$, $\mathcal{J}$ l'idéal qu'il engendre dans
$A \mathbin{\hat{\otimes}}_{k} A$, $\widetilde{\mathcal{J}}$ l'idéalisateur de
ce dernier ; la distinction typographique est de l'éditeur.

\pagerange{1}{1}
\section*{I. Le dictionnaire (page 1)}

La page énonce une équivalence de données. D'un côté :

\begin{itemize}
\item une $\otimes$-catégorie $C$ sur $k$, et un $\otimes$-foncteur
$k$-linéaire exact et fidèle $F' : C \to \operatorname{Mod}_{f}(k')$.
\end{itemize}

De l'autre :

\begin{enumerate}
\item[1\textsuperscript{o})] un schéma affine $G' = \operatorname{Spec} U$ sur
$k'$ muni d'une section, c'est-à-dire une $k'$-algèbre commutative $U$ munie
d'une augmentation $\varepsilon : U \to k'$ ;
\item[2\textsuperscript{o})] sur le pro-objet $A = \check{U}$, une
multiplication associative admettant $\varepsilon$ pour
unité\footnote{L'adjectif que la page accole ici est porté sous une rature et
la transcription le lit, en le signalant comme douteux, « commutative et
associative ». La commutativité est exclue par l'exemple du dossier lui-même
deux feuillets plus loin, où $A = \operatorname{End}_{k}(k')$ est une algèbre
de matrices ; et elle rendrait vide la difficulté que la page résout à sa
dernière ligne. La lecture retient donc l'associativité et l'unité, qui sont
ce dont la suite se sert.} ;
\end{enumerate}

soumis à deux conditions de multiplicativité, qui font l'objet de la section
suivante.

Ce qui est dit là, sans être dit, est ceci : $C$ est la catégorie des
$A$-modules de dimension finie sur lesquels $A$ opère continûment, $F'$ est le
foncteur d'oubli qui ne retient que le $k'$-espace vectoriel sous-jacent, et
c'est la multiplication de $A$ — l'article 2\textsuperscript{o}) — qui porte
toute cette partie-là de la structure. La comultiplication de $U$, dont la
multiplication de $A$ est la transposée, est la loi de composition du futur
groupoïde ; l'augmentation $\varepsilon$ en est l'unité\footnote{Cette moitié
du dictionnaire — une catégorie $k$-linéaire équivaut à une algèbre
linéairement topologisée ayant assez de représentations de dimension finie, ses
objets étant les modules permis — est construite en toutes lettres de sa main
ailleurs dans le fonds, aux pages 156 à 179 du dossier 10. Les feuillets d'ici
la supposent acquise et n'en refont rien.}.

Le produit tensoriel de $C$, lui, n'est pas dans la multiplication de $A$ : il
est dans la \emph{co}multiplication
\[
\Delta_{A} : A \longrightarrow A \mathbin{\hat{\otimes}} A,
\]
qui est la transposée de la multiplication de $U$ — laquelle est donnée
gratuitement par 1\textsuperscript{o}), puisque $U$ est une algèbre. C'est
pourquoi la page n'a pas à la postuler : elle a seulement à demander qu'elle
soit compatible au reste, et c'est la condition (ii).

En langage d'aujourd'hui, la donnée $(U, \varepsilon, A)$ est un
\emph{algébroïde de Hopf} — ou, du côté géométrique, un groupoïde affine
au-dessus de $\operatorname{Spec} k'$, agissant transitivement sur
$\operatorname{Spec} k'$ — moins son antipode, qui est précisément ce que les
pages 5 à 9 vont chercher\footnote{Les noms sont de nous. Ni « groupoïde », ni
« algébroïde de Hopf », ni « antipode » ne figurent sur ces feuillets, qui ne
citent d'ailleurs personne. Le traitement publié du cas où le foncteur fibre
n'existe que sur une extension est celui de Deligne, \emph{Catégories
tannakiennes} (1990), postérieur de plusieurs années à la datation du dossier ;
le cas neutre, où $k' = k$ et où le groupoïde est un groupe, est chez Saavedra
Rivano (1972) et Deligne--Milne (1982).}.

\pagerange{1}{1}
\section*{II. Les deux conditions, et pourquoi $A \mathbin{\hat{\otimes}}_{k'} A$ ne convient pas (page 1)}

\emph{Condition (i).} L'application $k' \to A$, $\lambda \mapsto
\lambda\varepsilon$, est multiplicative : l'unité de $A$ est $\varepsilon$, et
$k'$ s'envoie dans $A$ par un morphisme d'algèbres\footnote{Une note
interlinéaire, écrite d'une main plus rapide et en partie illisible, précise
que la multiplicativité s'entend « pour la structure de pro-anneau de $A$ ». La
lecture la reproduit dans ce sens.}.

\emph{Condition (ii).} La comultiplication $\Delta_{A}$, transposée de la
multiplication de $U$, est multiplicative — c'est-à-dire un morphisme
d'algèbres. C'est l'axiome de bigèbre, et c'est lui qui fait du produit
tensoriel de deux $A$-modules un $A$-module.

Reste à savoir dans quoi $\Delta_{A}$ prend ses valeurs, et c'est ici que la
page fait le travail intéressant. Elle écrit d'abord la flèche comme allant
dans $A \mathbin{\hat{\otimes}}_{k'} A$, puis se corrige aussitôt : les valeurs
sont dans
\[
\widetilde{\mathcal{J}}/\mathcal{J} \;\subset\;
\bigl(A \mathbin{\hat{\otimes}}_{k} A\bigr)\big/\mathcal{J},
\qquad
\mathcal{J} = J\cdot\bigl(A \mathbin{\hat{\otimes}}_{k} A\bigr),
\quad
J = \operatorname{Ker}(k' \otimes_{k} k' \to k'),
\]
et l'application $A \to \widetilde{\mathcal{J}}/\mathcal{J}$ est un morphisme
d'anneaux.

La correction est nécessaire, et voici pourquoi\footnote{La page ne donne pas
cette raison ; elle écrit les deux formules l'une après l'autre. L'explication
est de nous, mais la formule est d'elle, et c'est la formule qui est juste.}.
$A$ n'est pas commutative et $k'$ n'y est pas central : le corps opère sur $A$
de deux manières, celle qui vient de la structure de $k'$-algèbre de $U$ et
celle qui vient de la seconde patte du groupoïde, et ces deux actions ne
coïncident pas. Or le produit tensoriel $B \otimes_{R} B$ de deux copies d'un
anneau $B$ au-dessus d'un anneau $R$ n'est un anneau que si $R$ est central
dans $B$. Ici il ne l'est pas, et $A \mathbin{\hat{\otimes}}_{k'} A$ — qui
n'est autre que le quotient $(A \mathbin{\hat{\otimes}}_{k} A)/\mathcal{J}$ —
n'est qu'un module. Passer à l'idéalisateur $\widetilde{\mathcal{J}}$, c'est
prendre le plus grand sous-anneau de $A \mathbin{\hat{\otimes}}_{k} A$ dans
lequel $\mathcal{J}$ soit un idéal bilatère ; le quotient
$\widetilde{\mathcal{J}}/\mathcal{J}$ est alors un anneau, et c'est le seul
endroit où l'énoncé « $\Delta_{A}$ est multiplicative » ait un sens.

C'est, au dual près, le \emph{produit de Takeuchi} des algébroïdes de
Hopf\footnote{Takeuchi, \emph{Groups of algebras over $A \otimes \bar{A}$}
(1977), introduit pour la même raison le sous-anneau
$B \times_{R} B \subset B \otimes_{R} B$ formé des tenseurs qui commutent aux
deux actions de $R$. Le nom est de nous ; la page construit l'objet et ne le
nomme pas.}. La même difficulté reparaît, vue de l'autre côté, dans la note
marginale de la page 3 : la transposée $U \to U \otimes_{k} U$ de la
multiplication de $A$ n'y respecte pas les unités, et il le note
lui-même — « $\Delta_{U}(1) \neq 1$ en général ».

\pagerange{3}{3}
\section*{III. L'exemple : la descente de $k'$ à $k$ (page 3)}

Le plus petit cas est celui où l'on ne gagne rien à changer de base :
\[
C = \operatorname{Mod}_{f}(k),
\qquad
F' : M \longmapsto M \otimes_{k} k'.
\]
Le dictionnaire donne alors
\[
U = k' \otimes_{k} k',
\qquad
A = \check{U} = \operatorname{Hom}_{k'}(k' \otimes_{k} k', k')
\;\simeq\; k' \otimes_{k} \check{k}'
\;\simeq\; \operatorname{End}_{k}(k'),
\]
avec pour augmentation la multiplication $k' \otimes_{k} k' \to k'$.

Géométriquement, $\operatorname{Spec}(k' \otimes_{k} k')$ est le
\emph{groupoïde des paires} de $\operatorname{Spec} k'$ au-dessus de
$\operatorname{Spec} k$ : deux pattes, qui sont les deux projections, et une
composition qui recolle. Ses représentations sont exactement les données de
descente, et la catégorie de ces données est $\operatorname{Mod}_{f}(k)$ : on
retrouve $C$, ce qui est la descente galoisienne\footnote{Le nom est de nous.
La page se contente du calcul et de la remarque que les structures d'anneaux y
sont « les habituelles ».}.

Du côté de l'algèbre le même fait se lit autrement, et plus brutalement :
$\operatorname{End}_{k}(k')$ est une algèbre de matrices $M_{n}(k)$,
$n = [k' : k]$, donc Morita-équivalente à $k$, et ses modules de dimension
finie sont les $k$-espaces vectoriels de dimension finie. C'est bien $C$ qu'on
récupère, et on le récupère sans hypothèse de séparabilité : l'isomorphisme
$\operatorname{End}_{k}(k') \simeq M_{n}(k)$ ne demande que la finitude.

Ce petit exemple porte déjà tout ce que la page 1 avait de surprenant. $A$ y
est non commutative dès que $[k':k] > 1$, et $k'$ n'y est pas central : il s'y
plonge de deux manières, par la source et par le but, exactement les deux
actions qui rendaient nécessaire l'idéalisateur.

\pagerange{3}{3}
\section*{IV. La question du dossier (page 3)}

Un long trait sépare le calcul de ce qui suit, et ce qui suit est ouvert par un
crochet. Il tient en deux morceaux.

\emph{Premier morceau, technique.} On demande que le foncteur de changement de
base le long de la diagonale,
\[
M' \longmapsto M' \otimes_{k' \otimes_{k} k'} k',
\]
soit exact et fidèle sur les $A_{k'}$-modules de type fini, où
$A_{k'} = A \otimes_{k} k'$ est une algèbre sur $k' \otimes_{k} k'$. Il note
lui-même, entre parenthèses : « OK si $k'/k$ étale ».

C'est juste pour l'exactitude, et pour la raison suivante : $k'/k$ séparable
équivaut à $k' \otimes_{k} k'$ réduite, donc étale sur $k'$, donc produit de
corps ; $k'$ y est alors un facteur direct, donc plat, et le changement de base
est exact. Pour $k'/k$ non séparable, $k' \otimes_{k} k'$ est locale à
nilpotents et l'exactitude tombe\footnote{La page écrit « étale » là où
« séparable » suffirait : les deux conditions sont la même ici, la première
portant sur $k' \otimes_{k} k'$ et la seconde sur $k'/k$.}. La fidélité, en
revanche, est fausse telle quelle : dès que $k' \otimes_{k} k' \simeq k' \times
R$ avec $R \neq 0$, un module concentré sur le facteur $R$ est tué par le
foncteur. Ce que la condition demande réellement, c'est donc que les modules
provenant de $C$ soient concentrés sur la composante diagonale — et c'est
précisément le genre d'énoncé que la question suivante cherche à obtenir de $C$
seule\footnote{La distinction est de nous : la page demande « exact et
fidèle » sans restreindre la classe de modules.}.

\emph{Second morceau, de fond.} Quelles conditions \emph{intrinsèques} sur $C$
— l'insertion interlinéaire dit « intrinsèques », et c'est le mot qui compte —
garantissent tout cela ? Et il écrit la candidate, suivie de trois points
d'interrogation :
\begin{quote}
Suffit-il que $C$ soit \textbf{rigide}, et que $k \overset{\sim}{\to}
\operatorname{End}_{C}(1_{C})$ ? ? ?
\end{quote}
puis, doublement souligné, en dernière ligne du feuillet :
\begin{quote}
\textbf{Exprimer la rigidité de $C$ !}
\end{quote}

La réponse, telle que la théorie l'a donnée depuis, est : non, pas en
général — et il manque exactement une condition numérique. Une catégorie
tensorielle rigide, abélienne, $k$-linéaire, avec $\operatorname{End}(1) = k$,
admet un foncteur fibre sur une extension de $k$ si et seulement si, en
caractéristique nulle, la dimension de chacun de ses objets est un entier
$\geq 0$ ; et elle est alors la catégorie des représentations d'un groupoïde
affine du type que la page 1 décrit. En caractéristique positive la question
reste ouverte\footnote{C'est le théorème de Deligne, \emph{Catégories
tannakiennes} (1990), § 7 pour le critère de dimension et § 1 pour la
reconstruction par un groupoïde ; les contre-exemples sont les catégories
d'interpolation, où les dimensions ne sont pas entières. La réponse est de
nous, elle est postérieure au dossier, et rien sur ces feuillets ne l'annonce :
ils posent la question et la laissent ouverte.}.

Les deux pages qui suivent ne traitent pas de cette question-là. Elles
traitent de l'autre moitié du programme : \emph{étant} donnée la rigidité, que
devient-elle dans l'algèbre $A$ ?

\pagerange{5}{5}
\section*{V. La rigidité représentée : l'objet dualisant est $U$ (page 5)}

Le feuillet s'ouvre sur un titre de sa main, souligné : \emph{Rigidité}.

Soit $M \mapsto \check{M} = D(M)$ la dualité de $C$. Elle est compatible au
foncteur fibre, c'est une équivalence de catégories — donc exacte —, et il en
tire qu'elle est continue à droite, c'est-à-dire qu'elle transforme les limites
inductives en limites projectives. Un tel foncteur contravariant est
représentable : il existe un objet $I$ tel que
\[
D(M) \;\simeq\; \operatorname{Hom}_{A}(M, I).
\]
Reste à identifier $I$, et il suffit pour cela d'évaluer en $M = A$ :
\[
I \;=\; D(A) \;\simeq\; \check{A}_{g} \;=\; U.
\]
L'objet dualisant est donc $U$ lui-même — le dual du module régulier, que le
dictionnaire de la page 1 avait placé du côté « fonctions » et qui revient ici
du côté « objets »\footnote{La page porte, sous le signe d'isomorphisme, la
mention « compatibilité avec $F'$ », qui est ce qui autorise l'identification :
c'est parce que la dualité commute au foncteur fibre que le représentant se
calcule sur le module régulier.}.

Mais $U$ n'est pour l'instant qu'un module à droite : l'opération de $A$ qu'on
y lit vient, par transposition, de la multiplication à droite de $A$ sur
elle-même. Pour que $\operatorname{Hom}_{A}(M, U)$ soit encore un objet de $C$,
il faut que $U$ porte en plus une opération à gauche, commutant à celle-là. Il
faut donc faire de $U$ un \emph{bimodule}\footnote{La page écrit « une
structure de \emph{bimodule} module sur $A \otimes A$ », en biffant le premier
mot. L'algèbre qui classe les bimodules est $A \otimes_{k} A^{\circ}$, avec
l'opposée en second facteur, et c'est bien elle qui intervient : la ligne
suivante de la même page le montre, puisqu'elle produit $A^{\circ}$ et non $A$.
La lecture rétablit l'opposée.}.

\pagerange{5}{5}
\section*{VI. Ce que cette structure vaut : un anti-isomorphisme de $A$ (page 5)}

La structure supplémentaire cherchée se calcule d'un mot, parce que les
endomorphismes du module régulier sont connus :
\[
\operatorname{End}_{A}(U) \;\simeq\;
\operatorname{End}_{A}(A_{d})^{\circ} \;\simeq\; A^{\circ}.
\]
Se donner l'opération à gauche revient donc à se donner un morphisme
\[
A \longrightarrow A^{\circ},
\qquad a \longmapsto \check{a},
\]
c'est-à-dire un \emph{anti}-morphisme de $A$ dans elle-même : une application
qui renverse les produits : l'image de $ab$ est $\check{b}\,\check{a}$. Comme la
dualité est une équivalence et que son carré est isomorphe à l'identité, cet
anti-morphisme est bijectif et involutif\footnote{La page écrit « donc il faut
un donné $A \to A^{\circ}$ », avec « $a : a \mapsto \check{a}$ » au-dessus de
la flèche, et n'en dit pas plus ; l'involutivité, qui vient de $D^{2} \simeq
\operatorname{id}$, est ajoutée ici.}.

C'est l'\emph{antipode}. Dans le dictionnaire, la multiplication de $A$ est la
composition du groupoïde, la comultiplication est le produit tensoriel, et
cette involution qui renverse les produits est le passage à l'inverse : c'est
ce que la rigidité de $C$ devient une fois traduite. Le mot est de nous ; la
chose est sur la page\footnote{Pour une bigèbre de dimension finie l'énoncé
correspondant est classique : la catégorie de ses comodules de dimension finie
est rigide si et seulement si la bigèbre admet un antipode, c'est-à-dire est
une algèbre de Hopf. Le dossier démontre une moitié de cet énoncé — rigidité
$\Rightarrow$ antipode — dans le cadre plus large où la base est $k'$ et non
$k$ ; il n'énonce pas la réciproque.}.

On peut vérifier l'énoncé sur l'exemple de la page 3, ce que le dossier ne fait
pas. Là, le groupoïde est celui des paires et son inversion échange les deux
projections, donc échange les deux facteurs de $U = k' \otimes_{k} k'$ ; la
transposée de cet échange est un anti-automorphisme de $A =
\operatorname{End}_{k}(k')$, c'est-à-dire, par Skolem--Noether, la transposition
composée avec un automorphisme intérieur. L'inverse du groupoïde des paires est
bien la transposition des matrices.

\pagerange{7}{7}
\section*{VII. Le calcul du dual : $F'(D(M)) = \check{M}^{(k')}$ (page 7)}

Le feuillet suivant vérifie que la dualité ainsi obtenue est la bonne,
c'est-à-dire que le foncteur fibre l'envoie sur la dualité des espaces
vectoriels. Le calcul est une adjonction, écrite en trois lignes :
\begin{align*}
D(M) &\;=\; \operatorname{Hom}_{A}(M, U_{g})
\;=\; \operatorname{Hom}_{A}\bigl(M_{g},\,
\operatorname{Hom}_{k'}(A_{g}, k')\bigr) \\
&\;\simeq\; \operatorname{Hom}_{k'}\bigl(A_{d} \otimes_{A} M,\, k'\bigr).
\end{align*}
La page écrit le terme du milieu sous sa forme explicite : les formes
$k'$-bilinéaires $\varphi : A_{g} \times M_{g} \to k'$ telles que
\[
\varphi(\lambda u, x) = \lambda\,\varphi(u,x),
\qquad
\varphi(u\cdot v, x) = \varphi(u, v.x),
\]
$\lambda$ dans $k'$, $u$ dans $A$, $x$ dans $M$ — ce qui est exactement la
description des éléments du $\operatorname{Hom}$ adjoint.

Il reste à contracter, et c'est immédiat :
\[
A_{d} \otimes_{A} M \;\simeq\; M,
\qquad
x \longmapsto 1 \otimes x,
\qquad
\lambda(1 \otimes x) = \lambda \otimes x = 1 \otimes \lambda x,
\]
la seconde ligne étant la vérification que l'isomorphisme est $k'$-linéaire
pour la structure de $k'$-module venue de $A_{d}$. D'où la conclusion du
feuillet :
\[
F'\bigl(D(M)\bigr) \;=\; \check{M}^{(k')}.
\]
Le foncteur fibre envoie le dual catégorique sur le dual linéaire : c'est
l'énoncé, en une formule, qu'il est un $\otimes$-foncteur \emph{rigide}. Et
c'est l'anti-morphisme $a \mapsto \check{a}$ de la page précédente qui munit
$\check{M}^{(k')}$ de sa structure de $A$-module — la seule manière de faire
opérer $A$ sur un dual étant de renverser l'ordre des produits\footnote{La page
le dit en une demi-ligne dont deux mots sont illisibles ; le sens n'est pas en
doute, puisque c'est à cela que servait l'antipode.}.

\pagerange{7}{9}
\section*{VIII. La compatibilité qui reste, et l'arrêt (pages 7 et 9)}

Une chose n'est pas encore acquise. On sait maintenant que $F'(D(M))$ est le
dual linéaire de $F'(M)$, donc qu'il existe un accouplement d'évaluation
\[
F'(M) \otimes_{k'} F'\bigl(D(M)\bigr) \longrightarrow k' = F'(1_{C}).
\]
Mais rien ne dit encore que cet accouplement \emph{provienne} d'un morphisme de
$C$
\[
M \otimes D(M) \longrightarrow 1_{C},
\]
et sans cela la dualité obtenue n'est pas la dualité de la catégorie : elle
n'en est que l'ombre après application du foncteur fibre. La condition à
vérifier est que l'évaluation soit $A$-linéaire pour l'opération diagonale,
celle que donne la comultiplication $\Delta_{A}$ de la page 1.

La page 9 pose le problème sous cette forme et le réduit aussitôt au cas du
module régulier : il suffit de le prouver pour $M = A$, c'est-à-dire pour
l'accouplement canonique
\[
A^{k'} \otimes_{k'} U \longrightarrow k'.
\]
Puis la ligne suivante commence, « $A \otimes$ », et s'arrête. Le reste du
feuillet est blanc, et c'est le dernier feuillet de mathématiques du
dossier\footnote{La réduction à $M = A$ est licite parce que les deux membres
sont additifs en $M$ et que tout objet est quotient de sommes de copies du
module régulier ; la page ne le justifie pas.}.

L'équation qui restait à écrire porte, aujourd'hui, un nom : c'est l'axiome de
l'antipode,
\[
\textstyle\sum \check{a}_{(1)}\,a_{(2)} \;=\; \varepsilon(a),
\]
qui dit exactement que l'évaluation est un morphisme de comodules. Le dossier
s'arrête au seuil de cette vérification\footnote{L'écriture en notation de
Sweedler et le nom d'axiome de l'antipode sont de nous ; la page n'écrit ni
l'une ni l'autre. Ce qu'elle écrit, l'accouplement canonique sur le module
régulier, est le cas particulier dont cet axiome est la forme générale.}.

\pagerange{3}{9}
\section*{IX. Ce que le dossier laisse ouvert}

Quatre choses, et il vaut mieux les dire que les laisser deviner.

\emph{Les conditions intrinsèques.} C'est l'objet déclaré de la recherche au
bas de la page 3 — « on cherche conditions intrinsèques sur $C$ » — et aucune
n'est trouvée. Les feuillets font le chemin inverse : ils supposent la rigidité
et regardent ce qu'elle donne.

\emph{Les trois points d'interrogation.} La suffisance de « rigide, et $k
\overset{\sim}{\to} \operatorname{End}_{C}(1_{C})$ » est posée et laissée en
suspens. La réponse moderne, rappelée plus haut, est négative sans une
condition de dimension, en caractéristique nulle ; en caractéristique positive
elle n'est pas connue.

\emph{La dernière vérification.} Elle est réduite au module régulier et
abandonnée en cours de ligne.

\emph{La réciproque.} Que l'existence d'un antipode redonne la rigidité — la
moitié facile et la seule qui ferait du critère un critère — n'est ni énoncée
ni utilisée. Le programme du titre, « exprimer la rigidité de $C$ », est donc
tenu dans un sens seulement : la rigidité est traduite, elle n'est pas
caractérisée.

\end{document}
